% figure1_transport
\caption{First manifest-ordered fixed-rank V2 examples, HI1-prompt1-fixed and HT1-row1-fixed; selection did not use appearance. Images show unchanged saved pixels enlarged with nearest-neighbor interpolation and correct aspect ratios. A transports the canonical 16 $\times$ 16 grayscale array (256 bytes), not arbitrary original image-file bytes. Its literal leading text excerpt has display line wrapping only; the complete carrier is ../v2\_review/cases/HI1-prompt1-fixed/carrier.txt. B transports all 63 literal UTF-8 bytes through a 32 $\times$ 31 RGB8 PNG; the native model canvas is 32 $\times$ 32, with the shared 1 $\times$ 32 RGB row withheld. Fresh receivers used only their own saved carrier, pinned profile/model, shared prompt1 or row1, and retained encryption key. Sources and their byte comparison were evaluator-only. Row1 is 96 shared context bytes, not transmitted payload. The full examples and exact source/recovery paths are linked in the index.}
\label{fig:figure1_transport}

% figure2_recovery_rate
\caption{All six V2 cells, 20 held-out payload groups per direction. Recovery whiskers are the accepted two-sided 95\% Clopper--Pearson intervals (Bernoulli working model); rate whiskers are the accepted 2,000-draw, stratified payload-group percentile intervals. These preserve class/length strata and paired methods. Points show cell means. Goodput is 8 $\times$ source bytes divided by every delivered symbol only when the complete carrier authenticates and exactly recovers; otherwise it is zero. All 20 A1 text failures remain in the denominator at zero goodput. Recovered packet prefixes are not useful delivery (Figure S1). Text axes use bits/token; PNG axes use bits/channel value, not bits/pixel (multiply the PNG rates by three). Zero-width rate intervals reflect identical observed rates, not universal reliability. The canonical-image and fixed-slot formats constrain these small-sample results.}
\label{fig:figure2_recovery_rate}

% figure3_detectability
\caption{Each point is a saved V2 carrier or its matched ordinary control (20 of each per panel). Light connectors preserve payload-group pairing; horizontal offsets follow frozen group order only. Square open markers are ordinary controls; filled/colored symbols are stego, with crosses for failed A1 text transmissions. Scores are whole-carrier mean eligible-distribution surprisal under an observer who knows the exact model and shared context, without the encryption key. Each group owns one independent control trace; text methods reuse length-matched prefixes and PNG methods share the full ordinary PNG. Thus 120 method-specific links are only 40 independent traces across the study. No duplicate or unscorable V2 artifacts occurred. Annotations are the accepted higher-is-stego AUCs and 95\% grouped intervals, never flipped when below 0.5. All failed delivered A1 text carriers remain scorable. Panel y-scales differ. Complete empirical separation in the small fixed/gated text sample, including its collapsed bootstrap interval, does not establish perfect population detection. AUC near 0.5 does not establish security against this or other observers.}
\label{fig:figure3_detectability}

% figure4_gpu_performance
\caption{Development engineering benchmark, not V2: T1/T4/T5 contain 32/96/128 source bytes, all under row1. Fixed rank has three alternating-order matched timing repetitions per payload (nine comparisons). Colors identify execution mode and shapes repetitions; connectors in A/B join matched packets. Headline 5.133$\times$ speedup is the ratio of mean charged reference and graph encode-plus-decode time on the pinned RTX 5000 Ada / PyTorch 2.5.1+cu124 configuration. Charged fresh processes include imports, model hashing/loading, graph setup, CPU bookkeeping, serialization and teardown. Throughput is exactly recovered source bytes divided by both charged processes; the headline pools bytes and time across all nine repetitions, while C shows individual ratios. These are three development payloads, not nine independent research observations. All 11 benchmark comparisons (including single-payload gated/A1 in Table 4) had identical exact ID/order/probability stream hashes, delivered PNG bytes and recovered source bytes; all 22 recoveries passed. D shows maxima across fixed jobs separately for allocated tensors, reserved allocator storage and sampled process memory; they overlap and must not be added. Process memory is sampled, not a guaranteed true peak. No broad statistical speed claim or equivalence on other stacks is established; reference execution remains available.}
\label{fig:figure4_gpu_performance}

% figureS1_arithmetic_capacity
\caption{All 20 V2 A1 text and 20 A1 PNG cases, using retained audit observations every 128 consumed symbols and the final recorded endpoint. Dots are measured checkpoints; connecting segments are guides, not measured intermediate values. Axes have different carrier units and caps. Text stops at 2,048 tokens with 581--2,216 of the required 2,336 bits and never reaches zero-extension/termination; packet prefixes are not authenticated source delivery. PNGs reach the packet target before completing the remaining full 2,976-channel carrier. Curves stop at the packet endpoint; they do not extrapolate into ordinary completion. Useful packet-prefix bits are capped at 2,336; emitted zero suffix bits beyond that are kept separately in the supporting endpoint CSV and Table S1. Eligible surprisal and entropy checkpoint values are also in the data CSV, not invented from line segments. Isolated unchanged-interval steps in HI5/HI16 (failed text) and HT19 (successful PNG) were transient, longest run one; they do not establish sustained stagnation as the text-failure cause. Low accumulated information under the fixed cap remains the accepted explanation. A1 retains its separate, previously documented finite-precision stagnation limitation.}
\label{fig:figureS1_arithmetic_capacity}

% figureS2_text_consistency
\caption{V1 development comparison, separate from V2. Each column preserves one of 20 distinct canonical image payloads under both frozen prompts; 40 context-specific pairs are not 40 independent payloads. Within each pair, fixed-rank sequence and static arms transport the same immutable packet. Sequence eligibility checks the complete generated carrier prefix; static eligibility checks singleton tokens only, with the other filtering/coding rules held fixed. Static carriers were saved literally even when retokenization drifted; no repair, saved IDs or sender state reached receivers. Sequence recovered 20/20 under each prompt; static recovered 8/20 (forest) and 14/20 (soup), with 18 retained drift/replay failures. This motivates the tested main method’s saved-artifact consistency requirement, not a theorem that singleton filtering must always fail. The historical unavailable static matched-control prefix remains unavailable and is not scored or regenerated here.}
\label{fig:figureS2_text_consistency}
